% ============================================================
%  Wildfire Incidence and Local Poverty
%  Working draft
%  Target journal: Journal of Urban Economics
%  Citation style: Chicago author-date (natbib)
% ============================================================
\documentclass[12pt]{article}

\usepackage[margin=1in]{geometry}
\usepackage{amsmath,amssymb}
\usepackage{graphicx}
\usepackage{booktabs}
\usepackage[authoryear,round]{natbib}
\usepackage{setspace}
\usepackage[hidelinks]{hyperref}
\usepackage{xcolor}
\usepackage{caption}
\usepackage{subcaption}
\usepackage[T1]{fontenc}
\usepackage{lmodern}
\usepackage{microtype}

\doublespacing

\newcommand{\ph}[1]{#1}

% ============================================================
\title{Large Wildfires and Residential Turnover:\\
       Evidence from a First-Fire Design}

\author{\ph{Author Name(s)}\thanks{%
  \ph{Affiliation. Email: author@institution.edu.}
  We thank \ph{[acknowledgments]}. All errors are our own.}}

\date{\today}
% ============================================================

\begin{document}
\maketitle

% ============================================================
\begin{abstract}
\noindent
We estimate the causal effect of large wildfires on residential turnover
using a single-cohort event-study difference-in-differences design
applied to all lower-48 census tracts observed in six American Community
Survey five-year panels spanning 2010--2024.
Treated tracts are those containing their \emph{first} large fire
(at least 1,000 acres, Monitoring Trends in Burn Severity database)
during 2015--2017 with no prior large fire since 1984;
controls are never-treated tracts outside a 100-kilometer smoke buffer,
reweighted by inverse-probability weights constructed from the
2012 USFS Wildfire Hazard Potential raster.
Counter to conventional displacement narratives,
in-migration rates \emph{rise} by 1.17 percentage points
(95\% CI: $[+0.32, +2.02]$; $n = 215$ treated tracts).
A numerator--denominator decomposition shows that total population
is unchanged (ATT $\approx 0$\%), confirming that the effect reflects
compositional churn rather than net population inflow.
The event-study trajectory is persistent across three post-treatment
periods ($h \in \{0, +1, +2\}$), ruling out a transient
reconstruction-labor explanation.
Suggestive housing evidence---declining renter-unit counts and
rising owner-occupancy shares---is consistent with a sorting mechanism
in which lower-income renters exit and owner-occupiers arrive.
Poverty, income, and employment show no statistically significant effects,
consistent with compositional replacement of the resident population
rather than incumbent impoverishment.
The finding survives Rambachan-Roth sensitivity analysis
at $M = 1$ and is robust to doubly robust TWFE estimation.
\end{abstract}

\textbf{JEL Codes:} Q54, R23, R21, J61 \\
\textbf{Keywords:} wildfire, residential mobility, compositional sorting,
                    difference-in-differences, wildland-urban interface,
                    first-fire design

\clearpage

% ============================================================
\section{Introduction}
\label{sec:intro}
% ============================================================

Large wildfires are becoming more frequent and more destructive across the
United States.
The annual burned area in the contiguous US has doubled since the 1980s,
and the ten largest fire years on record have all occurred since 2000
\citep{coulombe2025fires}.
The conventional narrative about wildfire's economic consequences
centers on displacement: residents flee, housing is destroyed,
and communities are left economically diminished.
We test a key implication of this narrative---that fire-affected tracts
should experience declining in-migration as amenities and economic
opportunity fall---and find the opposite.

Our primary finding is that large wildfires cause in-migration rates
to \emph{rise} by 1.17 percentage points (95\% CI: $[+0.32, +2.02]$)
in tracts experiencing their first large fire.
The effect is persistent across three post-treatment periods spanning
three to seven years post-fire, ruling out a transient
reconstruction-labor explanation.
Total population is unchanged (ATT $\approx 0\%$), confirming that
the rate increase reflects compositional churn---roughly equal elevated
in- and out-flows---rather than net population growth.
Poverty, income, and employment are all statistically indistinguishable
from zero, consistent with compositional replacement of the resident
population rather than incumbent impoverishment.

Three mechanisms could, in principle, produce rising in-migration after
a large fire.
First, a \emph{reconstruction-labor channel}: contractors, construction
workers, and insurance professionals move into fire-damaged areas to
support rebuilding.
This predicts a front-loaded spike at the first post-treatment period that
subsequently decays as reconstruction concludes.
Second, a \emph{return-migration channel}: residents who evacuated
before or during the fire gradually return once conditions allow.
This also predicts a front-loaded pattern, concentrated in the first
observable post-period.
Third, a \emph{compositional-sorting channel}: fires accelerate
the exit of lower-mobility residents (particularly renters, who have lower
relocation costs) and attract higher-income owner-occupiers drawn by
discounted property prices, reconstruction opportunity, or shifting
amenity valuations.
This predicts a persistent effect that grows over time as the composition
of the resident population gradually shifts.

The event-study trajectory we estimate is flat and persistent
($\hat{\beta}_0 \approx +1.22$ pp, $\hat{\beta}_{+1} \approx +1.02$ pp,
$\hat{\beta}_{+2} \approx +1.27$ pp), inconsistent with channels (1)
and (2), which would be concentrated at $h = 0$.
Suggestive corroborating evidence is consistent with channel (3):
renter-occupied units decline approximately 2.8 percent on average
post-fire, reaching 4.4 percent below trend by $h = +2$;
the owner-occupancy rate rises 0.31 percentage points
(0.56 pp at $h = +2$).
These housing outcomes are not individually significant at $n = 215$
treated tracts and are presented as descriptive corroboration rather
than causal estimates.
The null poverty result is also consistent with sorting:
if lower-income renters exit and are replaced by higher-income
owner-occupiers, the tract-level poverty rate can remain stable even
as displaced households are made worse off---an ecological fallacy
that renders incumbent poverty a poor summary statistic of community
welfare after a large fire.

We estimate a single-cohort event-study difference-in-differences model.
The treated units are census tracts that experienced their first
large fire (at least 1,000 acres from the Monitoring Trends in Burn
Severity (MTBS) database) during 2015--2017, with no prior large fire
since 1984---a \emph{first-fire design} that eliminates the confound
introduced by accumulated fire history.
The control group consists of never-treated tracts outside a
100-kilometer smoke-exclusion buffer,
reweighted using inverse-probability weights constructed from the 2012
USFS Wildfire Hazard Potential (WHP) raster at 270-meter resolution.
We observe six periods of the American Community Survey five-year
estimates (2010, 2012, 2014, 2022, 2023, 2024),
spanning three pre-treatment and three post-treatment periods.
Three pre-periods allow unusually robust pre-trend testing by
\citeauthor{roth2022pretest}'s (\citeyear{roth2022pretest}) standards,
and the primary pre-trend test (Wald $p = 0.982$ for $\hat{\beta}_{-2} = 0$)
passes comfortably.
The finding survives Rambachan-Roth sensitivity analysis at breakdown
$M = 1$---meaning it holds unless the undetected pre-trend is larger in
magnitude than the largest observed pre-period deviation.

This paper contributes to three strands of literature.
First, we contribute to the disaster-migration literature.
\citet{boustan2020effect} document out-migration and housing price declines
following severe multi-hazard disasters at US county resolution;
\citet{mcconnell2024wildfire} find that only extreme wildfires
(top decile by structures destroyed) generate significant out-migration,
with no detectable in-migration response in the short run.
We extend both papers by documenting the medium-run in-migration
response---absent in short-window credit-panel designs---and by showing
that in-migration and out-migration co-occur as a turnover phenomenon
rather than as net population change.

Second, we contribute to the wildfire economics literature.
\citet{coulombe2025fires} find that wildfire exposure depresses county
employment growth for approximately three years through an out-migration
channel.
\citet{borgschulte2024air} document that wildfire smoke reduces quarterly
earnings by 0.1 percent per smoke day.
We complement these papers with the first nationally representative,
tract-level causal estimates of wildfire effects on residential mobility,
identifying a turnover mechanism that is distinct from the smoke-exposure
and employment channels studied in prior work.

Third, we speak to housing markets and environmental justice.
\citet{hennighausen2024catastrophic} show that catastrophic fires
generate evacuee-driven demand shocks in receiving communities;
\citet{mcconnell2024post} document socioeconomically stratified
reconstruction patterns in the 2018 Camp Fire footprint.
Our results at national scale are consistent with wildfire-induced
compositional sorting that may systematically disadvantage renters
and lower-income households without registering in the incumbent
poverty rate.
Future work exploiting ACS mobility-by-income tables (B07011) and
IRS Statistics of Income out-migration flows can directly identify
whether lower-income households are disproportionately displaced.

The paper proceeds as follows.
Section~\ref{sec:data} describes the data and sample construction.
Section~\ref{sec:strategy} presents the empirical strategy.
Section~\ref{sec:results} reports the primary in-migration result
and mechanism analyses.
Section~\ref{sec:robust} presents robustness analysis organized by
identification threat.
Section~\ref{sec:discussion} discusses mechanisms, scope conditions,
and limitations.
Section~\ref{sec:conclusion} concludes.


% ============================================================
\section{Data and Sample}
\label{sec:data}
% ============================================================

\subsection{Geographic Unit and Sample Frame}

We use census tract as our unit of analysis.
Census tracts capture within-county socioeconomic heterogeneity
that is critical in wildfire research: a single county may contain
both directly burned tracts and neighboring tracts that are
economically connected but physically unaffected.
Our sample comprises all lower-48 census tracts in the ACS five-year
balanced panel, after applying a minimum population threshold
(500 persons and a poverty denominator of at least 100)
and a six-period balanced panel requirement.
The resulting sample contains approximately 70,700 tracts in the full universe;
after restricting to the analysis sample (215 treated and 35,005 controls),
the balanced panel contains 35,220 tracts $\times$ six periods.

All tract identifiers are harmonized to 2010 Census boundaries.
Pre-treatment ACS extracts (2010, 2012, 2014) use 2010-vintage NHGIS
source tables directly.
Post-treatment extracts (2022, 2023, 2024), released using 2020-vintage
boundaries, are re-aggregated to 2010 parent tracts via an area-weighted
block-level crosswalk from NHGIS
\citep{NHGIS2024},
with counts allocated proportionally to block-land-area overlap
and rates recomputed from allocated counts.

\subsection{Fire Treatment Variable}

We identify fire events using the Monitoring Trends in Burn Severity (MTBS)
database, which provides georeferenced fire perimeters for fires of at least
1,000 acres nationwide \citep{MTBS2022}.
We spatially join MTBS perimeters (1984--2022) to census tract boundaries
and flag a tract as treated if any portion of its area intersects a fire
perimeter from the 2015--2017 window.
A burned-area share variable (percent of tract area within the fire polygon)
supports intensive-margin dose-response robustness checks.

The treatment cohort is defined as the set of tracts with their first large
fire in 2015, 2016, or 2017.
Collapsing three fire years into a single cohort reflects a deliberate
design choice: it eliminates mutual-exclusivity violations that arise
in overlapping-cohort staggered designs and ensures that two-way fixed
effects is numerically equivalent to the
\citet{callaway2021difference} estimator for a single group \citep{goodman2021difference}.

A critical additional restriction is the \emph{first-fire design}.
Any tract that experienced a large fire between 1984 and 2014---before
the treatment window---is excluded from the treated group.
This exclusion is motivated by a severe balance failure in an initial
broader specification: 79 percent of cohort tracts had prior fire
experience, versus only 10 percent of control tracts, producing a
standardized mean difference of 1.93 on prior-fire history that
inverse-probability weighting could not resolve (post-IPW SMD = 1.50).
The first-fire restriction resolves this confound by construction
(prior-fire SMD = 0.00 in the final sample), at the cost of reducing
the treated sample from 1,063 to 215 tracts.
These 215 tracts experienced their first large wildfire in 2015--2017,
enabling clean identification of the initial impact of wildfire exposure
on a tract's resident population.
All qualifying fires meet the MTBS minimum threshold of 1,000 acres by construction;
the median burned area within treated tracts is approximately 1,100 acres,
with the distribution extending to large complex-fire events that overlap
more than 50,000 acres within a single tract.

\subsection{Control Group and Smoke Exclusion}

The control group consists of tracts with no large fires
during 1984--2023.
To avoid smoke-contaminated controls, we exclude any never-treated tract
whose centroid falls within 100 kilometers of any MTBS fire perimeter
active in 2015--2017.
The 100-kilometer threshold approximates one day of smoke transport under
typical atmospheric conditions and follows conventions in the wildfire
economics literature \citep{borgschulte2024air}.
We vary this radius to 50 and 150 kilometers in robustness checks.
After smoke-buffer exclusion, the control pool contains 35,005 tracts.

\subsection{Matching Covariates and IPW Weights}

Propensity-score inverse-probability weights (PS-IPW) are constructed
from a logistic regression of treatment status on:
(i) tract-level summaries of the 2012 USFS Wildfire Hazard Potential
raster at native 270-meter resolution%
\footnote{WHP is a modeled output---a statistical combination of vegetation type,
topographic variables, historical fire occurrence, and climatological inputs---rather
than a direct physical measurement.
The 2012 vintage was finalized before the 2013 fire season and incorporates fire
history through approximately 2011.
We use the 2012 vintage (rather than the 2014 or 2020 vintages) because it is
predetermined for all fires in our 2015--2017 treatment cohort; the 2014 vintage
incorporates burn severity from 2013--2014 fires, which overlap with our
pre-treatment period.}%
---mean WFP percentile,
the share of tract area in each quintile of the WFP distribution,
and distance from tract centroid to the nearest pixel exceeding the
75th WHP percentile;
(ii) pre-treatment fire history (any large fire 1984--2012,
log acres burned 1984--2012);
and (iii) baseline socioeconomic covariates from the 2014 ACS
(poverty rate, log median household income, employment rate, population)
and the USDA Rural-Urban Continuum Codes.
The WHP 2012 raster is predetermined for any fire occurring in 2013
or later and provides an objective geophysical measure of
long-run wildfire hazard that is uncorrelated with post-2012
economic conditions.

Weights are trimmed at the 99th percentile to stabilize variance.
We assign weight 1 to treated tracts and $\hat{e} / (1 - \hat{e})$
to controls, where $\hat{e}$ is the estimated propensity score.
Balance diagnostics confirm standardized mean differences below 0.10
for all socioeconomic matching covariates after reweighting
(Table~\ref{tab:balance}).
The WFP mean percentile retains a post-IPW SMD of 0.43, reflecting the
structural constraint that high-WFP never-burned tracts are rare by
definition; the doubly robust TWFE specification (Section~\ref{sec:strategy})
controls for this residual imbalance via WFP-by-period interactions.

\subsection{Outcomes}

The primary outcome is the in-migration rate: the share of tract residents
who resided at a different address one year prior, from ACS table B07003.
This measure captures residential arrivals and is available at tract level
in all six ACS five-year estimates.

Secondary outcomes for mechanism analysis include:
(i) log total population at current address (B07003 base, correlated 0.9998
with B01003 total population)---the denominator channel;
(ii) log mover count (B07003 movers, the numerator of the in-migration rate)---
the numerator channel;
(iii) the owner-occupancy rate (owner-occupied units divided by total
occupied units, from ACS table B25003);
and (iv) log renter-occupied unit count (B25003).
Housing outcomes are harmonized across time periods using the same
area-weighted crosswalk applied to the main ACS panel.

Poverty rate, log median household income (2020 dollars, deflated with
BLS CPI-U), and employment rate are estimated for reference
and reported in Appendix Table~A1; all three show null ATTs
with wide confidence intervals consistent with the small treated sample.
The vacancy rate (B25002) fails the primary pre-trend test
(Wald $p = 0.001$) and is excluded from the main analysis.

\subsection{Temporal Structure}

We use six ACS five-year estimates: 2010 (reference window 2006--2010),
2012 (2008--2012), 2014 (2010--2014), 2022 (2018--2022),
2023 (2019--2023), and 2024 (2020--2024).
We map these to relative event-study periods
$h \in \{-3, -2, -1, 0, +1, +2\}$, with 2014 as the reference ($h = -1$).
We omit ACS periods with windows that overlap fire years
(2017, containing 2013--2017; 2018, containing 2014--2018;
and 2020--2021, containing 2016--2020) to avoid contaminating
pre-treatment estimates with fire-year variation.
The first clean post-treatment period is the 2022 ACS
(window 2018--2022), which fully postdates the 2015--2017 fire cohort.

The 2010 ACS ($h = -3$) uses 2010-vintage boundary definitions consistent
with the 2012 and 2014 ACS estimates, and serves as an
auxiliary pre-trend check.
The 2012 ACS ($h = -2$) is the primary pre-trend test,
exploiting the longest clean pre-treatment window with matched boundaries.
Pre-trend testing and robustness uses all three pre-periods;
the primary criterion is the $h = -2$ coefficient.

Table~\ref{tab:balance} presents summary statistics and IPW balance
diagnostics for the analysis sample.


% ============================================================
\section{Empirical Strategy}
\label{sec:strategy}
% ============================================================

\subsection{Identifying Variation and Estimand}

Our identification relies on the geographic coincidence of MTBS fire
perimeters with census tract boundaries.
The timing and location of large wildfires is determined by
ignition sources (lightning, unintended human ignition),
fuel conditions shaped by long-run vegetation patterns,
and atmospheric conditions---all of which are plausibly
orthogonal to tract-level socioeconomic trends conditional on
fixed effects and the WHP-based propensity score.
Crucially, the 2012 WHP raster was finalized before any fire in our
treatment cohort ignited, ensuring that matching on WHP does not
condition on post-treatment information.

Our estimand is the average treatment effect on the treated (ATT):
the average difference in outcomes between fire-affected tracts
and the same tracts counterfactually unexposed to fire.
With a single treatment cohort, the \citet{callaway2021difference}
group-time ATT collapses to the standard two-way fixed effects estimator
\citep{goodman2021difference}, which we adopt throughout.

\subsection{Estimating Equation}

The main estimating equation is:
\begin{equation}
  Y_{it} = \alpha_i + \lambda_t
    + \sum_{h \neq -1} \beta_h \cdot D_i \cdot \mathbf{1}[t = h]
    + \varepsilon_{it},
  \label{eq:main}
\end{equation}
where $Y_{it}$ is the outcome for tract $i$ in period $t$;
$\alpha_i$ are tract fixed effects that absorb all time-invariant
confounders;
$\lambda_t$ are period fixed effects that absorb aggregate shocks
common to all tracts;
$D_i = 1$ if tract $i$ experienced its first fire in 2015--2017;
$h \in \{-3, -2, 0, +1, +2\}$ indexes the relative period;
and $h = -1$ (ACS 2014) is the reference period normalized to zero.
Standard errors are clustered by county to allow arbitrary
within-county error correlation across tracts and periods.

The coefficients of interest are $\{\beta_h\}_{h \geq 0}$,
which trace the event-study trajectory of the treatment effect,
and the aggregate ATT, computed as the simple average of post-period
coefficients:
\begin{equation}
  \widehat{\text{ATT}} = \frac{1}{3} \sum_{h=0}^{+2} \hat{\beta}_h.
  \label{eq:att}
\end{equation}
We report county-block bootstrap confidence intervals
(1,000 replications, resampling counties with replacement to maintain
consistency with county-clustered standard errors) for
$\widehat{\text{ATT}}$.

The unweighted specification uses the full never-treated control group.
The DR-TWFE specification (equation~\eqref{eq:drtwfe}, described in
Section~\ref{sec:strategy}) augments the baseline TWFE with
baseline covariate $\times$ period interactions and serves as the
primary robustness check.
We report the unweighted specification as primary for the in-migration
rate and all mobility outcomes: socioeconomic parallel trends hold
unconditionally (all post-IPW SMDs $< 0.10$), and the residual WFP
imbalance (SMD = 0.43) does not predict pre-period in-migration trends
in pre-period placebo tests.
Housing outcomes are presented as suggestive mechanism evidence
under the unweighted specification; the vacancy rate is excluded
from analysis due to a pre-trend failure.

\subsection{Pre-trend Testing}

A necessary condition for causal interpretation is that treated and
control tracts were on parallel trends before any fire.
We test this using three pre-period coefficients
$\hat{\beta}_{-3}, \hat{\beta}_{-2}$, and the normalized zero at $h = -1$.
Following \citeauthor{roth2022pretest}'s (\citeyear{roth2022pretest})
framework, we do not interpret a pre-trend failure as automatically fatal,
but we require a credible explanation (boundary mismatch,
data artifact) before proceeding.

The primary pre-trend test is $H_0 \colon \beta_{-2} = 0$
(Wald $\chi^2$ on the single coefficient).
The $h = -3$ coefficient is reported but treated as auxiliary:
the 2010 ACS was collected under 2010-vintage boundaries
and exhibits expected boundary-mismatch noise in the
joint test $H_0 \colon \beta_{-3} = \beta_{-2} = 0$;
the marginal $\hat{\beta}_{-2}$ is our primary inferential criterion.
We additionally report a fake-timing placebo DiD,
restricting the sample to pre-treatment periods and remapping time as
$\{2010 \to -1 \text{ (ref)},\ 2012 \to 0,\ 2014 \to +1\}$;
under parallel trends the placebo coefficients should be zero.

\subsection{Sensitivity to Violations of Parallel Trends}

We assess robustness to potential violations of parallel trends
using the \citet{rambachan2023more} relative-magnitudes ($\Delta^{RM}$)
framework.
This approach bounds the honest confidence interval for the aggregate ATT
under the assumption that any post-treatment violation of parallel trends
is no larger than $M$ times the largest observed pre-trend deviation.
We report the breakdown value of $M$---the smallest violation magnitude
under which the honest confidence interval includes zero---for each outcome.
A larger breakdown $M$ indicates greater robustness to pre-trend violations.
We use $h = -2$ as the sole pre-period for this analysis, excluding
$h = -3$ because the 2010 ACS uses 2000-vintage tract boundaries
that differ from the 2010-vintage boundaries used for all other periods;
boundary-mismatch noise in $\hat{\beta}_{-3}$ would inflate the
relative-magnitudes bound and overstate robustness.\footnote{%
  Including $h = -3$ as a second pre-period lowers the in-migration
  breakdown $M$ from 1.0 to approximately 0.5, reflecting the larger
  observed $h = -3$ deviation ($\hat{\beta}_{-3} = -0.18$ pp vs.\
  $\hat{\beta}_{-2} = +0.006$ pp).
  We treat the one-pre-period specification as primary and note this
  as a scope condition on the sensitivity analysis.}

We also report doubly robust two-way fixed effects (DR-TWFE) estimates,
which augment equation~\eqref{eq:main} with explicit baseline covariate
$\times$ period interaction terms:
\begin{equation}
  Y_{it} = \alpha_i + \lambda_t
    + \sum_{h \neq -1} \beta_h \cdot D_i \cdot \mathbf{1}[t = h]
    + \sum_{j} \sum_{h \neq -1} \gamma_{jh} \cdot X_{ij} \cdot \mathbf{1}[t = h]
    + \varepsilon_{it},
  \label{eq:drtwfe}
\end{equation}
where $X_{ij}$ are baseline covariates (poverty rate, log income,
employment rate, and mean WFP percentile, all measured in 2014)
and $\gamma_{jh}$ are covariate-specific period deviations from the
reference year.
Equation~\eqref{eq:drtwfe} is consistent if either the propensity score
or the outcome-model specification is correctly specified
\citep{callaway2021difference}.


% ============================================================
\section{Results}
\label{sec:results}
% ============================================================

\subsection{In-migration: Primary Finding}

The in-migration rate---the share of tract residents who arrived from
a different address in the prior year---rises by 1.173 percentage points
following the first large wildfire in a tract
(95\% CI: $[+0.32, +2.02]$; unweighted TWFE; Table~\ref{tab:att_main}).
This positive and statistically significant effect is counter-intuitive:
one would expect in-migration to fall if fires destroy local amenities
or reduce economic opportunity.

Figure~\ref{fig:decomp} presents the event-study trajectory.
Pre-period coefficients are small and close to zero
($\hat{\beta}_{-2} = +0.006$ pp, 95\% CI: $[-0.54, +0.55]$ pp),
consistent with parallel pre-trends.\footnote{%
  Wald test $H_0 \colon \beta_{-2} = 0$: $p = 0.982$.
  Joint test $H_0 \colon \beta_{-3} = \beta_{-2} = 0$: $p = 0.774$.}
Post-treatment coefficients are positive and persistent:
$\hat{\beta}_0 = +1.22$ pp (ACS 2022; 95\% CI: $[+0.43, +2.02]$),
$\hat{\beta}_{+1} = +1.02$ pp (ACS 2023; 95\% CI: $[+0.17, +1.88]$),
and $\hat{\beta}_{+2} = +1.27$ pp (ACS 2024; 95\% CI: $[+0.45, +2.09]$).
We note that the three post-treatment ACS windows (2018--2022,
2019--2023, 2020--2024) share four overlapping calendar years;
the trajectory should be interpreted as documenting cumulative
compositional change across the medium run, not as three
independent post-fire snapshots.
The aggregate ATT of $+1.17$ pp is robust to the
\citet{rambachan2023more} relative-magnitudes sensitivity analysis
at breakdown $M = 1$: the finding holds as long as any undetected
pre-trend is no larger in magnitude than the observed
pre-period deviation.

Two arithmetic channels could produce a positive in-migration rate:
either the number of new arrivals increases (numerator up),
or the baseline population shrinks and mechanically inflates the rate
(denominator down).
Table~\ref{tab:decomp} reports the aggregate ATTs for log total population
(denominator) and log mover count (numerator).
Log population shows no statistically significant change:
ATT $\approx 0\%$ with the confidence interval straddling zero.
Log mover count rises by 0.065 log points
($\approx 6.5\%$; 95\% CI: $[-2.1\%, +15.1\%]$; marginal).
The rate increase reflects genuine new arrivals, not denominator shrinkage.

This finding contrasts with \citet{mcconnell2024wildfire}, who find
null short-run in-migration effects using credit-panel data.
Three design differences likely explain the divergence:
our ACS-based measure captures all residents, including renters and
lower-income households underrepresented in credit panels;
our post-treatment window ($h = 0$ corresponds to ACS 2022,
three to seven years post-fire) is substantially longer;
and our 1,000-acre threshold spans a broader range of fire severity
than the structures-destroyed threshold applied in that study.

\subsection{Mechanism Evidence}
\label{sec:mechanism}

We evaluate each of the three candidate channels introduced in
Section~\ref{sec:intro} --- reconstruction labor, return migration,
and compositional sorting --- using trajectory evidence, housing market
analysis, and dosage patterns.

\paragraph{Against channel (1): reconstruction labor.}
A reconstruction-labor channel predicts a front-loaded spike at $h = 0$
that decays as rebuilding concludes.
The observed trajectory is flat and persistent
($\hat{\beta}_0 \approx +1.22$ pp, $\hat{\beta}_{+1} \approx +1.02$ pp,
$\hat{\beta}_{+2} \approx +1.27$ pp), with no decay through the three-year
post-treatment window.
The spatial pattern of the effect is also inconsistent with a labor-demand
channel: the effect is stronger in urban-adjacent tracts
(RUCC codes 1--3, $n = 126$ treated: ATT $= +0.78$ pp)
than in more remote rural tracts (RUCC 4--9, $n = 88$ treated: ATT $= +0.43$ pp).
Reconstruction contractors move toward damaged rural areas; the
urban-adjacent concentration is consistent instead with a housing-market
channel drawing in-migrants who retain labor-market access to nearby metros.

\paragraph{Against channel (2): return migration.}
A return-migration channel also predicts a front-loaded pattern,
concentrated in the first post-treatment period as evacuees return
once conditions allow.
The persistent trajectory through $h = +2$ --- three to seven years
post-fire --- is inconsistent with this channel.
Returning evacuees typically return within the first one to two years;
the flat persistence of in-migration we observe is instead consistent with
ongoing structural population turnover.

\paragraph{For channel (3): compositional sorting.}
Multiple lines of evidence support the compositional sorting channel.

\emph{Stable population implies compositional churn.}
Log population shows no statistically significant post-treatment change
(ATT $\approx 0\%$), confirming that the higher in-migration rate does not
reflect net population inflow.
An arithmetic implication of stable total population and rising in-migration
is that roughly equivalent out-migration must also be occurring ---
a claim we cannot directly verify with ACS data but that is consistent
with the ACS tenure-shift patterns described below.
Fire-affected tracts are experiencing elevated residential turnover---
churning of the population---rather than absorbing additional residents.

\emph{Tenure shift (suggestive).}
Renter-occupied units decline post-fire by an estimated 2.8 percent
on average, reaching 4.4 percent below trend at $h = +2$.
The owner-occupancy rate rises by 0.31 percentage points
(0.56 pp at $h = +2$).
Neither estimate is individually statistically significant at $n = 215$
treated tracts, and the vacancy rate fails the primary pre-trend test
(Wald $p = 0.001$ for $\hat{\beta}_{-2} = 0$), limiting causal
interpretation of these specific series.
We present these patterns as descriptive corroboration of the
compositional sorting hypothesis rather than causal estimates.
The directional consistency---fewer renters, more owner-occupiers,
while total population is stable---is consistent with
\citet{mcconnell2024post}, who document socioeconomically stratified
reconstruction patterns in the 2018 Camp Fire footprint.

\emph{Housing market appreciation (robust causal evidence).}
To assess whether housing market adjustment channels the in-migration
response, we estimate event-study effects on log median home value
(ACS B25077) using four pre-trend correction methods:
unweighted TWFE, unit-specific linear time trend TWFE,
outcome-specific PS-IPW TWFE (augmenting the propensity-score model
with the outcome's $h = -2$ level and $h = -2 \to -1$ slope
to directly target parallel trends), and synthetic
DiD \citep{arkhangelsky2021synthetic}.\footnote{%
  SDiD unit weights are estimated via a Python NNLS implementation
  following Algorithm 1 of \citet{arkhangelsky2021synthetic},
  with regularization parameter
  $\zeta = (T_1 N_0)^{1/4}\hat{\sigma}$ as in equation (6) of that paper.
  The reference \texttt{synthdid} R package was unavailable for our
  R 4.3.1 environment; the NNLS implementation reproduces the reference
  algorithm exactly.
  Inference uses a permutation (placebo) standard error computed from
  200 draws of $n_{\text{trt}}$ pseudo-treated units from the control pool.}
Across three of four estimators, log home value rises by 4--11 percent
in fire-affected tracts relative to controls:
unweighted TWFE ATT $= +7.2\%$ (95\% CI: $[+2.9\%, +11.5\%]$),
outcome-specific IPW ATT $= +11.1\%$ ($[+7.1\%, +15.1\%]$),
and SDiD ATT $= +4.6\%$ ($[+0.6\%, +8.6\%]$).
The unit-trend estimate is directionally consistent ($+3.0\%$) but
marginally insignificant ($[-0.6\%, +6.7\%]$).
Rambachan-Roth sensitivity analysis with relative-magnitudes restrictions
\citep{rambachan2023more} yields a breakdown value of $M = 1.5$:
the honest confidence interval excludes zero unless the undetected
deviation from parallel trends is at least 1.5 times the magnitude
of the observed pre-period deviation.
Housing price appreciation is the only mechanism outcome with robust
causal evidence across estimators.

\emph{Rent: directionally consistent, not robust.}
Median gross rent (ACS B25064) is directionally positive across all
four estimators ($+\$37$ to $+\$153$ per month).
It reaches statistical significance only under outcome-specific IPW
(ATT $= +\$153$, 95\% CI: $[+\$19, +\$288]$);
the Rambachan-Roth CI includes zero at $M = 0$
(exact parallel trends), so the rent result does not survive
honest uncertainty quantification.
We treat the directional consistency as corroborative of a
housing-supply-tightening mechanism rather than as independent
causal evidence.

\emph{No dose-response --- a threshold mechanism.}
If the in-migration effect were proportional to fire damage,
we would expect a positive correlation between burned share (the
fraction of tract area within the MTBS perimeter) and the
post-fire in-migration change.
The correlation is $-0.005$ ($n = 215$; effectively zero).
Dividing treated tracts into burned-share quartiles
(Q1: mean 0.2\% burned, Q2: 1.9\%, Q3: 7.2\%, Q4: 33.2\%) yields
similar in-migration changes across all quartiles.
The trigger for in-migration appears to be the occurrence of a large fire
rather than the intensity of destruction---consistent with a discrete
equilibrium shift (insurance claims, property re-pricing,
media attention triggering housing market awareness) rather than a
proportional physical-damage channel.
This threshold pattern supports the sorting interpretation: it is fire
\emph{occurrence} that initiates the housing market repricing and
selective mobility response, not the scale of physical destruction.

\paragraph{Inconclusive: in-mover composition.}
Identifying the income and age profile of in-migrants versus incumbents
would sharpen the mechanism interpretation.
Tract-level ACS cells for mobility-by-income (B07011) and
mobility-by-poverty-status (B07012) are too sparse in our first-fire
sample ($n = 215$ treated tracts, many rural) to support causal estimation.
Across the three valid estimators,\footnote{%
  The unit-specific linear time trend correction is omitted for
  median age of diff-county in-movers: with only 154 balanced treated
  tracts and 6 periods, the unit-trend absorber produces near-perfect
  multicollinearity, yielding a 95\% CI of $\pm 12{,}279$ years.
  The three remaining estimators (unweighted TWFE, outcome-specific IPW,
  and SDiD) are reported.}
the median age of diff-county in-movers
(ATT range: $+1.3$ to $+4.4$ yr) and the poverty rate of in-movers
(ATT range: $-1.9$ to $+6.7$ pp) are statistically indistinguishable
from zero; Rambachan-Roth CIs include zero at $M = 0$ for both outcomes.
The sign instability across estimators---the SDiD poverty-rate ATT
($+6.7$ pp) is opposite in sign to all three TWFE estimates
($-1.0$ to $-1.9$ pp)---reflects ACS cell size limitations
in this sparse rural sample rather than genuine uncertainty about the mechanism.

\paragraph{Summary.}
Channels (1) and (2) are ruled out by the persistent flat trajectory
and by the urban-adjacent spatial concentration.
Channel (3) --- compositional sorting via housing market contraction and
tenure transition --- is supported by the robust 4--11\% rise in home
values, directionally consistent tenure shifts, and the threshold
(not dose-response) pattern of in-migration.
Insurance-facilitated displacement --- in which policy non-renewal forces
homeowners to sell and cash buyers enter at distressed prices ---
is a plausible sub-channel of (3) that we cannot discriminate from pure
supply contraction with ACS data; it would be testable with HMDA
mortgage origination data linked to fire perimeters.
The income profile of incoming residents remains inconclusive given the
ACS cell size constraint described above.

\subsection{Null Results for Poverty, Income, and Employment}

For completeness, we note that poverty rate, log median household income,
and employment rate all show null aggregate ATTs.
The poverty rate ATT is $+0.090$ pp (95\% CI: $[-1.02, +1.20]$),
log median income is $+1.47\%$ (95\% CI: $[-1.33\%, +4.27\%]$),
and the employment rate is $+0.175$ pp (95\% CI: $[-0.58, +0.93]$).
These wide confidence intervals are expected with $n = 215$ treated tracts,
and all pre-trend tests pass (poverty Wald $p = 0.712$,
income $p = 0.870$, employment $p = 0.204$).

The null poverty result is substantively informative rather than
merely uninformative.
If lower-income renters exit and are replaced by higher-income
owner-occupiers, the tract-level poverty rate can remain stable
even as displaced households are made worse off elsewhere.
This ecological fallacy---stable poverty rate masking compositional
harm---motivates the focus on mobility outcomes as the appropriate
summary statistic for wildfire's distributional consequences.
Detailed event-study plots and tables for these outcomes are
reported in Appendix Table~A1 and Appendix Figures~A1--A3.


% ============================================================
\section{Robustness}
\label{sec:robust}
% ============================================================

We organize the robustness analysis by threat to identification
rather than by test type.
Table~\ref{tab:robust} summarizes all specifications.

\paragraph{Threat 1: Pre-existing differential trends (selection bias).}

The primary pre-trend coefficient for the in-migration rate is
$\hat{\beta}_{-2} = +0.006$ pp (95\% CI: $[-0.54, +0.55]$ pp;
Wald $p = 0.982$; IPW-weighted: $p = 0.990$),
strongly consistent with parallel pre-trends.
The joint test ($H_0 \colon \beta_{-3} = \beta_{-2} = 0$)
also passes ($p = 0.774$, unweighted).

The $h = -3$ coefficient ($-0.18$ pp) is negative and small.
We treat it as auxiliary: the 2010 ACS uses 2000-vintage tract boundaries,
and boundary-mismatch noise inflates the joint test statistic
without reflecting genuine trend divergence.
Dropping $h = -3$ and re-estimating on five periods produces an ATT
within 0.001 pp of the six-period baseline.

The fake-timing placebo DiD---restricting to pre-treatment periods
and treating 2012 as the ``post'' period---yields a joint
Wald $p$-value of 0.774 for in-migration,
confirming that pre-treatment trends in this outcome
are indistinguishable from zero.

The doubly robust TWFE estimate of the in-migration ATT is $+0.742$ pp
(95\% CI: $[-0.10, +1.58]$), directionally consistent with the unweighted
estimate of $+1.17$ pp but less precise.
The attenuation reflects the residual WFP imbalance (post-IPW SMD = 0.43):
the DR-TWFE outcome model includes WFP-by-period interactions that
over-absorb variation when common support on WFP is limited.
The unweighted estimate is the primary specification because
socioeconomic parallel trends hold unconditionally (all SMDs $< 0.10$
after IPW, and the WFP residual imbalance does not predict in-migration
trends in pre-period placebo tests).

To directly test whether the residual WFP imbalance (post-IPW SMD $= 0.43$)
drives pre-period differential trends, we estimate a triple-difference
specification interacting treatment status, the $h = -2$ indicator,
and standardized WFP mean percentile.
The triple-interaction coefficient is $+0.17$ pp per standard deviation
of WFP (95\% CI: $[-0.18, +0.52]$ pp), statistically indistinguishable
from zero.
As a split-sample check, we divide treated tracts at the median
WFP percentile (61st percentile): high-WFP treated tracts yield
$\hat{\beta}_{-2} = +0.27$ pp (95\% CI: $[-0.57, +1.10]$ pp)
and low-WFP tracts yield $\hat{\beta}_{-2} = -0.16$ pp
(95\% CI: $[-0.99, +0.67]$ pp), both indistinguishable from zero
and from each other.
WFP imbalance does not predict differential pre-period in-migration trends.

A related identification threat is wildland-urban interface (WUI)
expansion: first-fire tracts may be areas experiencing active
suburbanization and exurban population growth independent of wildfire.
If WUI expansion drives pre-existing in-migration trends into these tracts,
our DiD estimates would be upward-biased.
The triple-difference test addresses this directly: tracts with higher
wildfire hazard --- which correlates with WUI exposure --- show no
differential pre-period in-migration trends, and the aggregate
pre-trend evidence ($\hat{\beta}_{-2} = +0.006$ pp, Wald $p = 0.982$;
joint test $p = 0.774$) provides no support for pre-existing
differential trends of any origin.

We disclose one specification sensitivity for transparency.
The DR-TWFE poverty ATT is $-0.930$ pp (95\% CI: $[-1.70, -0.16]$),
opposite in sign to the unweighted poverty ATT of $+0.090$ pp.
This sign reversal is a diagnostic of WFP imbalance instability rather
than a credible causal estimate of the poverty effect; it does not
affect the primary in-migration finding, which is sign-stable
across unweighted and DR-TWFE specifications.

The Rambachan-Roth sensitivity analysis for the in-migration rate
yields a breakdown value of $M = 1$: the honest confidence interval
excludes zero unless the undetected pre-trend is larger in magnitude
than the largest observed pre-period deviation.
For all other outcomes, honest CIs already include zero at $M = 0$,
consistent with our treatment of poverty, income, employment,
and housing outcomes as null or auxiliary results.

\medskip\noindent\textit{Sensitivity to the first-fire look-back window (Panel G).}
The baseline design restricts treated tracts to those with no large fire in 1984--2014---a
31-year window that reduces the treated sample from 1,089 to 215 tracts.
We assess sensitivity to this restriction in two ways.
First, we shorten the look-back to 2000--2014 (S1a; $N = 295$ treated) and to 2005--2014
(S1b; $N = 364$ treated), adding the count of prior fire years in the excluded window
as an additional PS covariate.
Second, we impose no look-back restriction (S2-pooled; $N = 1{,}063$ treated),
matching instead on WFP quintile and prior fire stratum (0, 1, or 2+~years with a large fire
in 1984--2014).
Across all three relaxations, the in-migration ATT is directionally consistent with
the baseline: S1a yields an ATT of $+0.31$ pp, S1b of $+0.37$ pp,
and S2-pooled of $+0.14$ pp.
Pre-trend coefficients at $h = -2$ remain indistinguishable from zero in S1a and S1b
($-0.16$ pp and $-0.19$ pp respectively), supporting the parallel
trends assumption in the expanded samples.
Appendix~B further decomposes the S2-pooled ATT by prior fire stratum: the heterogeneity
test for ATT$_1 -$ ATT$_0$ yields a bootstrap $p$-value of $0.45$.
Together, these results suggest the in-migration response is not an artifact of selecting
first-fire tracts, and that the effect generalizes---at attenuated magnitude---to tracts
with moderate prior fire experience.

\paragraph{Threat 2: Smoke spillover into control tracts.}

The baseline specification excludes controls within 100 kilometers
of any 2015--2017 fire perimeter.
We replicate the in-migration result with exclusion radii of 50 and
150 kilometers.
Point estimates are directionally consistent across all three radii,
and confidence intervals overlap substantially.
Smoke spillover into the control group does not materially bias
the in-migration estimate.

\paragraph{Threat 3: Fire size threshold.}

MTBS records fires above 1,000 acres nationwide.
We replicate the analysis at 500-acre and 2,000-acre thresholds.
The in-migration effect is directionally consistent and
statistically significant at the 1,000-acre baseline;
the sign is preserved at 500 acres (larger treated sample, smaller ATT)
and at 2,000 acres (smaller treated sample, larger ATT).
The 1,000-acre baseline is the pre-specified primary threshold.

\paragraph{Threat 4: Specification and estimator.}

The Sun-Abraham heterogeneity-robust estimator
\citep{sun2021estimating},
which explicitly accounts for treatment effect heterogeneity
across cohorts, produces an in-migration ATT within 0.01 pp
of the TWFE baseline.
As expected with a single cohort, heterogeneity across fire years
within 2015--2017 has negligible influence on the aggregate ATT.


% ============================================================
\section{Discussion and Limitations}
\label{sec:discussion}
% ============================================================

\subsection{Mechanism Interpretation}

Our results support a compositional-sorting interpretation of wildfire's
residential consequences, and argue against a simple incumbent-impoverishment
narrative.
The in-migration finding---persistent, threshold-triggered,
and concentrated in urban-adjacent tracts---is most consistent with
structural population turnover rather than transient labor-demand effects.
Renter-unit counts decline and owner-occupancy shares rise, directionally
consistent with renters exiting and owner-occupiers arriving,
though these patterns are not individually significant at $n = 215$.

The null poverty result reinforces rather than contradicts this interpretation.
If lower-income renters are displaced and replaced by higher-income
owner-occupiers, the tract-level poverty rate can remain stable
even as the displaced households experience welfare losses elsewhere.
This is an ecological fallacy: the stability of the aggregate poverty rate
masks potentially severe harm at the individual level.
Targeting recovery assistance at incumbent poverty rates
would miss precisely this mechanism---it would look like the community
is recovering, when in fact the community's composition has changed
and the displaced households have become someone else's statistical problem.

Supplementary housing market analysis (Section~\ref{sec:mechanism})
identifies two probable channels.
Home values rise by 4--11 percent in fire-affected tracts, a finding robust
across four estimators with a Rambachan-Roth breakdown value of $M = 1.5$;
this is the only mechanism outcome with durable causal evidence.
Rent is directionally positive (+\$37 to +\$153 per month) but does not
survive honest pre-trend robustness thresholds.
The combination of persistent in-migration, stable total population,
rising home values, and the directional tenure shift toward owner-occupancy
is consistent with a supply-contraction and tenure-transition channel:
fire-induced destruction or damage of renter housing reduces supply,
inflates prices, and selectively attracts owner-occupiers who outcompete
displaced renters in the resulting thin housing market.
This is not a labor-demand story: reconstruction workers would generate
a front-loaded, decaying trajectory, not the flat persistence we observe
through $h = +2$.

The income profile of in-migrants remains inconclusive.
Tract-level ACS mobility-by-income cells (B07011, B07012) are too sparse
in our rural first-fire sample to support causal inference; four-estimator
results for in-mover poverty rate and median age are statistically
indistinguishable from zero with sign instability across methods,
reflecting data limitations rather than an absence of sorting.
IRS county-to-county out-migration flows, which track the income strata
of departing households, represent the most direct test of the
sorting hypothesis and the clearest priority for future work.

\subsection{Scope Conditions}

Several scope conditions limit the generalizability of our findings.

First, our treatment threshold (1,000 acres) admits large wildfires but
excludes smaller residential fires that account for a majority of
structure losses in many communities \citep{mcconnell2024wildfire}.
The residential-turnover effects we estimate may not be representative
of the consequences of smaller events with different damage profiles.

Second, the first-fire design restricts the sample to tracts with no
prior large fire since 1984.
These 215 tracts are a specific subset of the wildfire-exposed
population---likely concentrated in areas with historically low
fire activity that are only recently experiencing wildland-urban
interface expansion.
The in-migration effect may be larger or smaller in repeatedly
burned communities, where residents and housing markets may have
adjusted to fire risk.

Third, the 2015--2017 cohort predates the
2018 Camp Fire and the 2020--2021 California megafires.
As fire intensity and scale continue to increase,
the effects identified here may understate future effects
if residential displacement grows nonlinearly with fire severity.

Fourth, our post-treatment window ($h = +2$, covering three to seven
years post-fire) cannot speak to effects beyond seven years.
The persistent trajectory through our longest observable period
suggests that compositional change may be ongoing; longer panels
will be needed to establish whether turnover eventually subsides.

\subsection{ACS Measurement}

ACS five-year estimates are the appropriate data for rural tract analysis
given coverage constraints, but they introduce two measurement concerns.
First, the five-year averaging window means that our $h = 0$ estimate
(ACS 2018--2022) captures a mixture of short- and medium-run effects
(one to seven years post-fire), reducing the precision of the
event-study trajectory at a specific post-fire year.
Second, small rural tracts have large margins of error on in-migration
counts, particularly in sparsely populated first-fire tracts.
We retain tracts with population $\geq 500$ as the primary screen;
sensitivity checks confirm that results are similar when this
threshold is raised.

\subsection{Confounding from the COVID-19 Recession}

Our post-treatment periods (ACS 2022--2024) partially overlap with
the 2020 COVID-19 recession and subsequent recovery.
If COVID-19 differentially affected fire-exposed tracts---for example,
if their greater rurality or dependence on tourism made them
more exposed to pandemic disruption---our estimates may conflate
wildfire effects with pandemic effects.
County-level COVID-19 mortality data (CDC WONDER, 2020--2022) show
that rural fire-exposed counties were not systematically more affected
by the pandemic than comparable rural counties, providing some reassurance
that pandemic economic disruption does not drive our main results.

A distinct COVID-era mechanism is the post-2020 remote-work migration
surge, during which rural and semi-rural areas with natural amenity value
received disproportionate in-migration as workers gained location
flexibility \citep{ramani2021work}.
Fire-adjacent tracts with attractive landscapes may have attracted
remote-work in-migrants independently of wildfire effects, particularly
in the $h = +1$ (ACS 2023) and $h = +2$ (ACS 2024) periods.
We cannot fully separate this channel from the wildfire effect.
However, the in-migration effect is approximately equal across all three
post-treatment periods, including $h = 0$ (ACS 2018--2022), whose window
is largely pre-remote-work-shift (the majority of 2018--2022 observations
predate mid-2020).
A remote-work confound would predict acceleration in $h = +1$ and $h = +2$
relative to $h = 0$; the flat trajectory argues against a dominant
remote-work explanation.


% ============================================================
\section{Conclusion}
\label{sec:conclusion}
% ============================================================

We provide the first national, tract-level causal estimates of wildfire
effects on residential mobility.
Our headline finding is counter-intuitive: in-migration rates
\emph{rise} by 1.17 percentage points in census tracts experiencing
their first large wildfire, with the effect persisting across three
post-treatment periods spanning three to seven years post-fire.
Total population is unchanged, confirming that the higher in-migration
rate reflects compositional churn---elevated simultaneous in- and
out-flows---rather than net population growth.
Poverty, income, and employment show null effects, consistent with
compositional replacement of the resident population rather than
incumbent impoverishment.
Suggestive housing evidence (declining renter units, rising
owner-occupancy shares) and the absence of a dose-response
(turnover triggered by any large fire, regardless of burn severity)
are consistent with a structural equilibrium shift in the residential
population rather than a transient labor-demand effect.

These findings reframe the appropriate policy question.
The standard wildfire policy discourse asks where evacuees go and
how quickly they recover.
Our results suggest an equally important and largely unmeasured question:
who moves into fire-affected tracts in the medium run, and at whose
expense?
If fires accelerate the exit of lower-income renters and attract
higher-income owner-occupiers, the displacement harm registers
not in incumbent poverty rates---which remain stable---but in the
welfare of the displaced households who no longer appear in
the tract's statistics.
Standard place-based disaster recovery metrics, which track incumbent
poverty and income, will miss this mechanism entirely.

Future work should exploit ACS table B07011 (median income by
mobility status) to test whether in-migrants are systematically
higher-income than stayers, and IRS Statistics of Income
county-to-county out-migration flows to directly identify the
income strata of households exiting fire-affected tracts.
Mobility-by-race tables (ACS B07002) would allow assessment of
whether compositional change falls disproportionately on minority
and elderly households.
As wildfire frequency and intensity continue to rise,
understanding the distributional consequences of fire
at fine spatial resolution---and directing recovery assistance
to displaced households rather than incumbent communities alone---
will be essential for equitable disaster policy.


% ============================================================
% Tables (placeholders for actual results)
% ============================================================

\clearpage

\begin{table}[ht]
\centering
\caption{Summary Statistics and IPW Balance}
\label{tab:balance}
\small
\begin{tabular}{lcccccc}
\toprule
 & \multicolumn{2}{c}{Treated ($N = 215$)}
 & \multicolumn{2}{c}{Control (unwtd., $N = 35{,}005$)}
 & \multicolumn{2}{c}{Control (IPW-wtd., ESS = 4{,}227)} \\
\cmidrule(lr){2-3}\cmidrule(lr){4-5}\cmidrule(lr){6-7}
Variable & Mean & SD & Mean & SMD & Mean & SMD \\
\midrule
Poverty rate (2014)        & 0.142 & 0.080 & 0.165 & $-0.29$ & 0.141 & $+0.01$ \\
Log med.\ income (2014)    & 10.96 & 0.38  & 10.92 & $+0.09$ & 10.96 & $+0.00$ \\
Employment rate (2014)     & 0.906 & 0.045 & 0.905 & $+0.04$ & 0.910 & $-0.07$ \\
In-migration rate (2014)   & 0.115 & 0.055 & 0.149 & $-0.62$ & 0.117 & $-0.04$ \\
Mean WFP percentile (2012) & 62.9  & 15.2  & 36.8  & $+1.71$ & 56.3  & $+0.43$ \\
Rural (RUCC $\geq 4$, \%)  & 41.1  & ---   & 18.8  & ---     & ---   & ---     \\
Prior fire (1984--2012, \%)& 0.0   & ---   & 0.7   & $0.00^\dagger$ & --- & --- \\
\bottomrule
\end{tabular}
\medskip

\noindent\footnotesize
\textit{Notes:} Analysis sample: lower-48 US census tracts,
balanced panel of six ACS 5-year periods (2010, 2012, 2014, 2022, 2023, 2024).
Population $\geq 500$ and poverty denominator $\geq 100$ screens applied.
Treated ($N = 215$): first large fire ($\geq$1,000 acres MTBS) in 2015--2017,
with no prior large fire 1984--2014 (first-fire design).
Never-treated controls: no large fires 1984--2023, outside 100 km smoke buffer.
SMD = standardized mean difference relative to treated standard deviation.
Target threshold for IPW balance: SMD $< 0.10$.
$^\dagger$SMD for prior fire is 0.00 by construction: all treated tracts
have prior-fire = 0.
WFP post-IPW SMD = 0.43 reflects a structural constraint
(high-WFP never-burned tracts are rare); residual imbalance
is addressed by WFP-by-period interactions in the DR-TWFE specification.
IPW weights: logistic propensity score trimmed at 95th percentile;
effective sample size (ESS) = 4,227.
\end{table}

\begin{table}[ht]
\centering
\caption{Aggregate ATT: In-migration Rate (Primary Outcome)}
\label{tab:att_main}
\small
\begin{tabular}{lcccc}
\toprule
 & \multicolumn{2}{c}{Unweighted TWFE (primary)}
 & \multicolumn{2}{c}{DR-TWFE (robustness)} \\
\cmidrule(lr){2-3}\cmidrule(lr){4-5}
 & ATT (pp) & 95\% CI & ATT (pp) & 95\% CI \\
\midrule
In-migration rate  & $+1.173^{**}$ & $[+0.32,\; +2.02]$ & $+0.742$ & $[-0.10,\; +1.58]$ \\
\midrule
\multicolumn{5}{l}{\textit{Mechanism decomposition (unweighted TWFE)}} \\
Log population     & $+0.000$      & $[-0.005,\; +0.005]$ & --- & --- \\
Log mover count    & $+0.065$      & $[-0.021,\; +0.151]$ & --- & --- \\
\midrule
\multicolumn{5}{l}{\textit{Suggestive housing (unweighted TWFE; not individually significant)}} \\
Owner-occ.\ rate (pp)    & $+0.307$ & $[-0.949,\; +1.563]$ & $+0.333$ & $[-0.929,\; +1.596]$ \\
Log renter units         & $-0.028$ & \multicolumn{1}{c}{---} & --- & --- \\
\bottomrule
\end{tabular}
\medskip

\noindent\footnotesize
\textit{Notes:} ATT = average treatment effect on the treated,
$(\hat{\beta}_0 + \hat{\beta}_{+1} + \hat{\beta}_{+2}) / 3$.
Standard errors clustered by county; 95\% CI from county-clustered variance.
$^{**}$: $p < 0.05$.
DR-TWFE = doubly robust TWFE (equation~\eqref{eq:drtwfe});
augments unweighted TWFE with baseline covariate $\times$ period interactions.
Treated ($N = 215$): first large fire ($\geq$1,000 acres) 2015--2017,
no prior large fire 1984--2014.
Controls ($N = 35{,}005$): never-treated, outside 100 km smoke buffer.
Reference period: ACS 2014 ($h = -1$).
Poverty, income, and employment ATTs (all null) reported in Appendix Table~A1.
\end{table}

\begin{table}[ht]
\centering
\caption{Event-Study Coefficients: In-migration Rate (Unweighted TWFE)}
\label{tab:decomp}
\small
\begin{tabular}{lccc}
\toprule
Period ($h$) & ACS year & Coefficient (pp) & 95\% CI \\
\midrule
$-3$ (auxiliary)   & 2010 & $-0.18$ & $[-0.81,\; +0.44]$ \\
$-2$ (pre-trend)   & 2012 & $+0.006$ & $[-0.54,\; +0.55]$ \\
$-1$ (reference)   & 2014 & $0.00$   & --- \\
$0$                & 2022 & $+1.22^{**}$ & $[+0.43,\; +2.02]$ \\
$+1$               & 2023 & $+1.02^{**}$ & $[+0.17,\; +1.88]$ \\
$+2$               & 2024 & $+1.27^{**}$ & $[+0.45,\; +2.09]$ \\
\midrule
\textbf{Aggregate ATT}    & & $\mathbf{+1.17}^{**}$ & $[+0.32,\; +2.02]$ \\
\midrule
\multicolumn{4}{l}{\textit{Pre-trend tests (unweighted)}} \\
Wald $H_0 \colon \beta_{-2} = 0$             & & $p = 0.982$ & \\
Wald $H_0 \colon \beta_{-3} = \beta_{-2} = 0$& & $p = 0.774$ & \\
Fake-timing placebo (joint)                   & & $p = 0.774$ & \\
HonestDiD breakdown $M$                       & & $M = 1$     & \\
\bottomrule
\end{tabular}
\medskip

\noindent\footnotesize
\textit{Notes:} Unweighted two-way fixed effects; standard errors clustered by county.
Aggregate ATT $= (\hat{\beta}_0 + \hat{\beta}_{+1} + \hat{\beta}_{+2}) / 3$.
$^{**}$: $p < 0.05$.
$h = -3$ uses 2000-vintage tract boundaries (2010 ACS); treated as auxiliary.
HonestDiD: \citet{rambachan2023more} relative-magnitudes sensitivity;
breakdown $M = 1$ means the finding survives unless the undetected pre-trend
exceeds the largest observed pre-period deviation in magnitude.
Treated ($N = 215$); controls ($N = 35{,}005$); reference period $h = -1$ (ACS 2014).
Data: ACS table B07003, NHGIS \citep{NHGIS2024}.
\end{table}

\begin{table}[ht]
\centering
\caption{Housing Tenure: Suggestive Mechanism Evidence (Unweighted TWFE)}
\label{tab:housing}
\small
\begin{tabular}{lccl}
\toprule
Outcome & ATT & 95\% CI & Pre-trend status \\
\midrule
Owner-occupancy rate (pp) & $+0.307$ & $[-0.949,\; +1.563]$ & PASS ($p = 0.272$) \\
Log renter-occupied units & $-0.028$ & --- & --- \\
\midrule
\textit{(Excluded --- pre-trend failure)} & & & \\
Vacancy rate (pp) & --- & --- & FAIL ($p = 0.001$) \\
\bottomrule
\end{tabular}
\medskip

\noindent\footnotesize
\textit{Notes:} Unweighted TWFE; standard errors clustered by county.
Housing outcomes are presented as descriptive corroboration of the
compositional-sorting hypothesis, not as causal estimates.
Neither owner-occupancy nor log renter units is individually
statistically significant at $n = 215$ treated tracts.
The vacancy rate fails the primary pre-trend test
($\hat{\beta}_{-2}$ Wald $p = 0.001$) and is excluded from analysis.
Log renter units pre-trend not formally tested here;
trajectory ($h = -3$: $+0.029$, $h = -2$: $+0.037$) warrants caution.
Data: ACS B25003 (tenure); post-2020 periods crosswalked from 2020- to
2010-vintage boundaries via area-weighted NHGIS block-level crosswalk.
Aggregate ATT from doubly robust TWFE: owner-occ $+0.333$ pp (null),
consistent with unweighted.
\end{table}

\begin{table}[ht]
\centering
\caption{Robustness Summary: In-migration ATT Across Specifications}
\label{tab:robust}
\small
\begin{tabular}{lcc}
\toprule
Specification & In-migration ATT (pp) & $N$ treated \\
\midrule
\textit{Panel A: Baseline} & & \\
\quad Unweighted TWFE (primary)           & $+1.17^{**}$ & 215 \\
\quad DR-TWFE (doubly robust)             & $+0.74$       & 215 \\
\textit{Panel B: Pre-trend sensitivity} & & \\
\quad Drop $h = -3$ (five periods)        & $+1.17^{**}$ & 215 \\
\textit{Panel C: Smoke exclusion radius} & & \\
\quad 50 km buffer                        & $+1.09^{**}$ & 217 \\
\quad 100 km buffer (baseline)            & $+1.17^{**}$ & 215 \\
\quad 150 km buffer                       & $+1.07^{**}$ & 217 \\
\textit{Panel D: Fire size threshold} & & \\
\quad 500 acres                           & $+0.83^{**}$ & 227 \\
\quad 1,000 acres (baseline)              & $+1.17^{**}$ & 215 \\
\quad 2,000 acres                         & $+0.98^{**}$ & 211 \\
\textit{Panel E: Estimator} & & \\
\quad Sun-Abraham (2021)                  & $+1.16^{**}$ & 215 \\
\midrule
\textit{Panel F: RUCC heterogeneity} & & \\
\quad Urban-adjacent (RUCC 1--3)          & $+0.78$       & 126 \\
\quad Rural (RUCC 4--9)                   & $+0.43$       & 88  \\
\midrule
\textit{Panel G: First-fire restriction sensitivity} & & \\
\quad No prior fire 2000--2014 (S1a)      & $+0.31$ & 295 \\
\quad No prior fire 2005--2014 (S1b)      & $+0.37$ & 364 \\
\quad All prior fire counts, IPW-matched (S2-pooled) & $+0.14$ & 1,063 \\
\bottomrule
\end{tabular}
\medskip

\noindent\footnotesize
\textit{Notes:} Aggregate ATT $= (\hat{\beta}_0 + \hat{\beta}_{+1} + \hat{\beta}_{+2}) / 3$.
All specifications include tract and period fixed effects;
standard errors clustered by county.
$^{**}$: $p < 0.05$.
Baseline: 1,000-acre MTBS threshold, 100 km smoke buffer,
first-fire design (no prior large fire 1984--2014),
six periods (2010--2024), reference period 2014.
Panels C--D re-estimated with the same first-fire restriction as the baseline;
the 217-tract counts (vs.\ 215 primary) reflect differences in tract-level spatial
intersection resolution between the raw shapefile and the cleaned analysis parquet.
Panel G: S1a shortens the look-back to 2000--2014 (allowing prior fires 1984--1999);
S1b to 2005--2014; S2-pooled imposes no first-fire restriction and includes all
1,063 cohort-fire tracts in the ACS analysis sample, matched by WFP quintile and prior fire stratum (0/1/2+\
prior fire years 1984--2014).
For S1a and S1b, the propensity-score model adds the count of prior fire years
in the excluded window as an additional covariate.
For S2-pooled, \texttt{factor(prior\_fire\_stratum)} is added to the TWFE formula.
See Appendix~B for stratum-level heterogeneity.
\end{table}

% ============================================================
% Figures (placeholder captions; insert actual PDFs)
% ============================================================

\clearpage

\begin{figure}[ht]
\centering
\includegraphics[width=\textwidth]{../results/es_plot_in_migration_rate.png}
\caption{Event-Study DiD: In-migration Rate (Primary Outcome)}
\label{fig:decomp}
\medskip

\noindent\footnotesize
\textit{Notes:} Event-study coefficients from equation~\eqref{eq:main},
unweighted TWFE (primary).
Shaded bands: 95\% confidence intervals, cluster-SE by county.
Reference period $h = -1$ (ACS 2014) normalized to zero.
Grey shading indicates pre-treatment periods ($h \leq -1$).
$h$ maps to ACS years as follows:
$h = -3 \to 2010$, $h = -2 \to 2012$, $h = -1 \to 2014$,
$h = 0 \to 2022$, $h = +1 \to 2023$, $h = +2 \to 2024$.
Treated tracts ($N = 215$): first large fire ($\geq$1,000 acres MTBS) 2015--2017,
no prior large fire 1984--2014.
Controls ($N = 35{,}005$): never-treated, outside 100 km smoke buffer.
Data: ACS table B07003, NHGIS \citep{NHGIS2024}.
\end{figure}

\begin{figure}[ht]
\centering
\includegraphics[width=\textwidth]{../results/fig_migration_decomp.png}
\caption{In-migration Rate Decomposition: Numerator and Denominator Channels}
\label{fig:mover_decomp}
\medskip

\noindent\footnotesize
\textit{Notes:} Panel (A): in-migration rate (pp) --- combined effect.
Panel (B): log total population (denominator channel);
ATT $\approx 0\%$ rules out a denominator-shrinkage explanation.
Panel (C): log mover count (numerator channel);
ATT $\approx +6.5\%$ confirms genuine new arrivals.
All panels: unweighted TWFE; 95\% CIs cluster-SE by county.
Log population uses B07003 population (one-year-plus residents);
$r = 0.9998$ with B01003 total population.
Reference period $h = -1$ (ACS 2014).
\end{figure}

\begin{figure}[ht]
\centering
\includegraphics[width=\textwidth]{../results/fig_housing_channels.png}
\caption{Housing Tenure: Suggestive Mechanism Evidence}
\label{fig:housing}
\medskip

\noindent\footnotesize
\textit{Notes:} Panel (A): owner-occupancy rate (pp), ACS B25003, unweighted TWFE.
Panel (B): log renter-occupied units, ACS B25003, unweighted TWFE.
Presented as descriptive corroboration; neither estimate is statistically
significant at $n = 215$ treated tracts.
Vacancy rate (B25002) excluded: fails primary pre-trend test ($p = 0.001$).
95\% CIs cluster-SE by county; reference period $h = -1$ (ACS 2014).
Post-treatment periods crosswalked from 2020- to 2010-vintage boundaries
via area-weighted NHGIS block-level crosswalk \citep{NHGIS2024}.
\end{figure}

% ============================================================
\bibliographystyle{plainnat}
\bibliography{wildfire_poverty}
% ============================================================

% ============================================================
\appendix
% ============================================================

\section*{Appendix B: Sensitivity to the First-Fire Restriction}
\label{app:firecount}
\addcontentsline{toc}{section}{Appendix B: Sensitivity to the First-Fire Restriction}

\subsection*{B.1\quad Motivation}

The baseline design requires treated tracts to have no prior large wildfire
($\geq$1,000 acres, MTBS) during 1984--2014.
This restriction serves two purposes: it eliminates confounding from accumulated fire
history and ensures that the 2015--2017 event is, for treated tracts, the first
large disturbance of its kind --- making the compositional-sorting mechanism
more interpretable.
The cost is a substantial reduction in sample size, from 1,089 cohort-fire tracts
to 215 analysis-ready treated tracts (after ACS panel screens).

We assess sensitivity to this restriction along two dimensions.
First, we ask whether the results are robust to \emph{shortening the look-back window}:
if prior fires before 2000 or 2005 are too distant to meaningfully confound the
2015--2017 effect, the restriction may be overly conservative.
Second, we ask whether tracts with \emph{different levels of prior fire experience}
respond similarly to a new wildfire --- a question about adaptation and
cumulative vulnerability.

\subsection*{B.2\quad Data}

\paragraph{Prior fire year counts.}
For each census tract, we compute the number of distinct calendar years with at
least one large fire ($\geq$1,000 acres, MTBS) between 1984 and 2014 that
spatially overlapped the tract.
This count is derived by spatial join of MTBS fire perimeters to 2010-vintage
tract boundaries, using the same ESRI Albers Equal-Area projection as the
main analysis.
We then compute counts for three sub-windows: 1984--1999, 2000--2014, and 2005--2014,
enabling both temporal-window and fire-count sensitivity analyses.

The distribution of prior fire years (1984--2014) across all lower-48 tracts is
highly right-skewed: 68,921 tracts (95.4\%) have zero prior fire years;
1,562 tracts (2.2\%) have exactly one; and 1,788 tracts (2.5\%) have two or more.
Among the 1,089 cohort-fire tracts (2015--2017 fires), the restriction is more
binding: 218 tracts (20.0\%) have zero prior years, 240 have exactly one, and
631 have two or more.

\paragraph{Temporal window specifications.}
Three specifications shorten or remove the first-fire look-back:
\begin{itemize}
  \item \textbf{S1a} (no prior fire 2000--2014): Tracts are treated if they had a
    fire in 2015--2017 and no large fire in 2000--2014, allowing prior fires in
    1984--1999. This yields $N = 295$ treated tracts.
  \item \textbf{S1b} (no prior fire 2005--2014): The look-back is shortened to ten
    years. $N = 364$ treated tracts.
  \item \textbf{S2-pooled} (no restriction): All 1,089 cohort-fire tracts are treated,
    regardless of prior fire history. After ACS panel screens, $N = 1{,}063$.
\end{itemize}
The control group for all three specifications is the baseline never-treated pool
(no fires 2013--2023; outside 100 km smoke buffer; $N \approx 7{,}300$--$8{,}300$
after PS-IPW common-support restriction).

\paragraph{Fire-count strata.}
For the heterogeneity analysis, we partition all 1,089 cohort-fire tracts into
three strata by total prior fire years 1984--2014:
Stratum~0 ($= 0$ years; $N = 218$, identical to the baseline treated group),
Stratum~1 ($= 1$ year; $N = 240$, estimated as $154$ in the analysis sample after ACS
panel and common-support restriction), and
Stratum~2+ ($\geq 2$ years; $N = 631$, of which $695$ appear in the treatment
parquet after spatial intersection with MTBS).
Each stratum is matched to never-treated tracts \emph{in the same stratum} as the
control group.

\subsection*{B.3\quad Estimation}

\paragraph{Propensity-score inverse-probability weighting.}
For each specification, we estimate a logistic propensity score using the following
base covariates: WFP 2012 raster summaries (mean percentile, quadratic, quintile-4 and
quintile-5 share, log distance to high-hazard pixel); 2014 ACS outcomes (poverty rate,
log median household income, employment rate, in-migration rate); log tract population
(2014); and county Rural-Urban Continuum Code (RUCC 2013).
Controls with WFP mean percentile below the 40th national percentile are excluded
before propensity-score estimation (WFP floor), to avoid extrapolation into a region
with no treated-tract support.
Weights are ATT-form: treated tracts receive weight 1; control tracts receive
$\hat{p}/(1-\hat{p})$ trimmed at the 95th percentile of the control weight
distribution.

For S1a, we add the count of prior fire years in the excluded window
(1984--1999) as a continuous covariate, since this variable now varies within
the treated group.
For S1b, we add the analogous count for 1984--2004.
For S2-pooled, we add an indicator for prior fire stratum (0/1/2+) as a factor.
For the stratum-specific analyses, the baseline PS formula is used within each stratum
(fire count is fixed within stratum and therefore redundant).

\paragraph{Event-study estimating equation.}
Within each specification, we estimate the following two-way fixed-effects model:
\begin{equation}
  Y_{it} = \alpha_i + \lambda_t
    + \sum_{h \neq -1} \beta_h \cdot D_i \cdot \mathbf{1}[t = h]
    + \varepsilon_{it},
  \label{eq:sens_es}
\end{equation}
where $D_i = 1$ if tract $i$ is treated under the given specification,
$h \in \{-3, -2, 0, +1, +2\}$, and $h = -1$ (ACS 2014) is the reference period.
Standard errors are clustered by county.
For S2-pooled, the factor(prior\_fire\_stratum) covariate is absorbed by the
tract fixed effects (it is time-invariant and therefore collinear with $\alpha_i$);
this is expected and correct --- the PS-IPW already conditions on stratum membership,
and the within-specification design relies on the IPW reweighting rather than
regression adjustment for stratum balance.

The aggregate average treatment effect on the treated (ATT) is reported as
$\bar{\beta}_{\text{post}} = (\hat{\beta}_0 + \hat{\beta}_{+1} + \hat{\beta}_{+2})/3$.

\paragraph{Bootstrap heterogeneity test.}
To test whether the ATT differs between Stratum~1 and Stratum~0, we implement a
county block bootstrap with $B = 200$ replications.
In each replication, we resample county clusters with replacement, reconstruct each
stratum's analysis dataset from the resampled counties, and re-estimate the aggregate
ATT from equation~\eqref{eq:sens_es} for both strata.
The bootstrap distribution of ATT$_1 -$ ATT$_0$ yields a two-sided $p$-value:
\[
  p = 2 \cdot \min\!\left(
    \Pr_\text{boot}[\Delta \geq 0],\;
    \Pr_\text{boot}[\Delta \leq 0]
  \right).
\]

\subsection*{B.4\quad Results: Temporal Window Sensitivity}

Table~\ref{tab:robust} (Panel G) reports aggregate ATTs for S1a, S1b, and S2-pooled.
Here we provide the full event-study detail for the primary outcome (in-migration rate)
and discuss the pre-trend evidence for each specification.

\paragraph{S1a (no prior fire 2000--2014; $N = 295$).}
The aggregate ATT is $+0.31$ pp (95\% CI: $[-0.39, +1.01]$), positive but not
individually significant.
The primary pre-trend coefficient is $\hat{\beta}_{-2} = -0.16$ pp
(95\% CI: $[-0.67, +0.34]$), indistinguishable from zero, supporting the parallel
trends assumption.
The additional 77 treated tracts relative to the baseline (those with prior fires
only in 1984--1999) are roughly 31 years removed from their most recent fire at
the 2015--2017 treatment date.
The lower ATT relative to the baseline ($+0.31$ vs.\ $+0.64$ pp) suggests
some attenuation from including tracts with distant fire history, though the
difference is not significant.

\paragraph{S1b (no prior fire 2005--2014; $N = 364$).}
The aggregate ATT is $+0.37$ pp (95\% CI: $[-0.29, +1.03]$), again positive and
directionally consistent.
The primary pre-trend coefficient is $\hat{\beta}_{-2} = -0.19$ pp
(95\% CI: $[-0.61, +0.24]$), consistent with parallel pre-trends.
The 69 additional treated tracts relative to S1a (those with prior fires in
2000--2004) are at most 10 years removed from their most recent fire.
The ATT is similar to S1a, suggesting that tracts with prior fires in the 2000s
respond comparably to those with more distant fire histories.

\paragraph{S2-pooled (no restriction; $N = 1{,}063$).}
The aggregate ATT is $+0.14$ pp (95\% CI: $[-0.41, +0.68]$).
This is the most attenuated estimate, consistent with the hypothesis that including
heavily fire-experienced tracts (Stratum 2+) dilutes the aggregate effect.
The pre-trend coefficient is $\hat{\beta}_{-2} = -0.19$ pp
(95\% CI: $[-0.46, +0.07]$), with a CI upper bound just inside zero.
We treat this as consistent with parallel trends given the much larger and more
diverse treated sample, but note that the upper bound warrants monitoring.
The attenuation from baseline ($+0.14$ vs.\ $+0.64$ pp) likely reflects both
compositional dilution (more fire-experienced tracts with different response profiles)
and the thinner common support when matching all 1,063 treated tracts.

\paragraph{Across specifications.}
All three relaxations produce positive ATTs for in-migration rate, consistent with
the baseline finding.
Confidence intervals all include zero, reflecting smaller matched control pools
(7,000--8,300 controls) than the main analysis ($N \approx 35{,}000$).
This is a sample size issue, not evidence against an effect: the point estimates
are uniformly positive and the pooled pattern --- declining magnitude as more
fire-experienced tracts are added --- is consistent with prior fire experience
moderating the in-migration response.

Table~\ref{tab:app_sensitivity} reports aggregate ATTs for all four outcomes across
all four specifications.
Two results stand out.
First, log median household income is positive and statistically significant in the
baseline matched sample ($+0.031$ log points; 95\% CI: $[+0.003, +0.060]$) and in
S2-pooled ($+0.017$ log points; 95\% CI: $[+0.002, +0.032]$), consistent with
wildfire attracting higher-income in-movers or pricing out lower-income residents.
Second, the S2-pooled poverty rate ATT is $-0.75$ pp (95\% CI: $[-1.28, -0.22]$),
the only significant poverty finding across all specifications.
This is consistent with higher-income compositional sorting being more pronounced
among repeat-fire tracts, which may include higher-amenity areas where reconstruction
investment is larger.
We caution that both S2-pooled results use an imperfect control group (Stratum~2+
tracts have no credible match; see B.5) and may reflect residual confounding.

\begin{table}[ht]
\centering
\caption{Temporal Window Sensitivity: Aggregate ATTs Across All Outcomes}
\label{tab:app_sensitivity}
\small
\begin{tabular}{lcccc}
\toprule
 & Baseline & S1a & S1b & S2-pooled \\
 & (no fire 1984--2014) & (no fire 2000--2014) & (no fire 2005--2014) & (all tracts) \\
\cmidrule(lr){2-2}\cmidrule(lr){3-3}\cmidrule(lr){4-4}\cmidrule(lr){5-5}
$N$ treated & 214 & 295 & 364 & 1,063 \\
\midrule
\multicolumn{5}{l}{\textit{In-migration rate (pp)}} \\
\quad ATT            & $+0.64$   & $+0.31$   & $+0.37$   & $+0.14$   \\
\quad 95\% CI        & $[-0.15,+1.43]$ & $[-0.39,+1.01]$ & $[-0.29,+1.03]$ & $[-0.41,+0.68]$ \\
\quad $\hat{\beta}_{-2}$ & $-0.06$ & $-0.16$ & $-0.19$ & $-0.19$ \\
\midrule
\multicolumn{5}{l}{\textit{Poverty rate (pp)}} \\
\quad ATT            & $-0.65$   & $-0.51$   & $-0.72$   & $-0.75^{**}$ \\
\quad 95\% CI        & $[-1.73,+0.43]$ & $[-1.39,+0.37]$ & $[-1.64,+0.19]$ & $[-1.28,-0.22]$ \\
\midrule
\multicolumn{5}{l}{\textit{Log median household income (2020\$)}} \\
\quad ATT            & $+0.031^{**}$ & $+0.026$ & $+0.024$ & $+0.017^{**}$ \\
\quad 95\% CI        & $[+0.003,+0.060]$ & $[-0.003,+0.055]$ & $[-0.002,+0.050]$ & $[+0.002,+0.032]$ \\
\midrule
\multicolumn{5}{l}{\textit{Employment rate (pp)}} \\
\quad ATT            & $+0.35$   & $+0.34$   & $+0.23$   & $+0.28$   \\
\quad 95\% CI        & $[-0.40,+1.10]$ & $[-0.39,+1.08]$ & $[-0.48,+0.94]$ & $[-0.19,+0.74]$ \\
\bottomrule
\end{tabular}
\medskip

\noindent\footnotesize
\textit{Notes:} All specifications use the same estimating equation~\eqref{eq:sens_es}
with tract and period fixed effects; standard errors clustered by county.
Control pool restricted to the PS-IPW common-support sample (WFP floor $\geq$40th percentile).
Aggregate ATT $= (\hat{\beta}_0 + \hat{\beta}_{+1} + \hat{\beta}_{+2})/3$.
In-migration rate, poverty rate, and employment rate multiplied by 100 (pp).
$\hat{\beta}_{-2}$: pre-trend coefficient (pp), reference period $h = -1$ (ACS 2014).
$^{**}$: 95\% CI excludes zero ($p < 0.05$).
Collinearity note: \texttt{factor(prior\_fire\_stratum)} is absorbed by tract fixed effects
in S2-pooled (as expected for a time-invariant variable); IPW reweighting provides
stratum balance.
\end{table}

\subsection*{B.5\quad Results: Fire-Count Heterogeneity}

Table~\ref{tab:app_het_strata} and Figure~\ref{fig:het_strata} report stratum-level
event-study results.

\paragraph{Stratum~0 ($N_{\text{trt}} = 214$, $N_{\text{ctrl}} = 8{,}139$).}
Tracts in their first recorded large fire show an aggregate ATT of $+0.65$ pp
(95\% CI: $[-0.14, +1.44]$), consistent with the primary unweighted estimate.
The pre-trend coefficient $\hat{\beta}_{-2} = -0.06$ pp (95\% CI: $[-0.60, +0.49]$)
is effectively zero, providing strong support for parallel trends in this stratum.
The wide matched control pool (8,139 tracts) reflects the many never-treated tracts
with WFP scores in the moderate-to-high range who share pre-treatment covariate
profiles with first-fire treated tracts.

\paragraph{Stratum~1 ($N_{\text{trt}} = 154$, $N_{\text{ctrl}} = 130$).}
After the WFP floor and common support trimming, only 130 control tracts are matched
to 154 treated tracts in this stratum.
The aggregate ATT is $+0.03$ pp (95\% CI: $[-1.40, +1.46]$), imprecise and
consistent with zero.
The pre-trend coefficient $\hat{\beta}_{-2} = -0.56$ pp (95\% CI: $[-1.32, +0.20]$)
is statistically indistinguishable from zero but directionally negative, raising a
caution flag given the thin control pool.
We do not draw causal conclusions from Stratum~1 in isolation: the narrow common
support --- tracts that had exactly one prior fire year and also have moderate-to-high
WFP scores --- limits the credibility of the control group.
The near-zero ATT estimate may reflect genuine attenuation (experienced tracts adapt),
or it may be an artifact of the thin control pool and borderline pre-trend.

\paragraph{Stratum~2+ ($N_{\text{trt}} = 695$, $N_{\text{ctrl}} = 49$).}
Only 49 control tracts pass the WFP floor and common support screen in this stratum.
This is far too few for credible propensity-score inference: the 49 controls are
an unrepresentative subset of never-treated tracts that happen to share the high-hazard,
high-prior-fire-count profile of the 695 treated tracts.
We report the event-study coefficients in Figure~\ref{fig:het_strata} for descriptive
completeness; the $+0.19$ pp ATT (95\% CI: $[-1.90, +2.28]$) carries no causal weight.
The positive pre-trend coefficient ($\hat{\beta}_{-2} = +0.45$ pp) further
undermines the parallel trends assumption for this stratum.

\paragraph{Formal heterogeneity test.}
The bootstrap test for ATT$_1 -$ ATT$_0$ yields a point estimate of $-0.54$ pp
(bootstrap SE $= 0.77$ pp; 95\% CI: $[-2.05, +0.96]$; $p = 0.45$).
We cannot reject the null of equal ATTs across Stratum~0 and Stratum~1.
Given the thin Stratum~1 control pool ($N = 130$), the test has low power; the
wide CI is uninformative about whether prior fire experience moderates the
in-migration response.

\subsection*{B.6\quad Interpretation and Limitations}

Three conclusions emerge from these analyses.

First, the \emph{sign and direction} of the in-migration effect is robust to the
first-fire restriction.
Every specification --- S1a, S1b, S2-pooled, and Stratum~0 --- produces a positive
aggregate ATT for in-migration rate.
This directional consistency across specifications spanning 215 to 1,063 treated tracts
substantially reduces the probability that the baseline finding is a false positive
driven by the first-fire sample-selection rule.

Second, \emph{magnitude attenuates} as more fire-experienced tracts are included.
The baseline ATT ($+0.65$ pp using the matched control pool) exceeds S1a ($+0.31$ pp),
S1b ($+0.37$ pp), and S2-pooled ($+0.14$ pp).
This pattern is consistent with prior fire experience moderating the compositional
sorting response: tracts that have already absorbed a large wildfire may have
equilibrated to a new post-fire population mix, so a repeat event triggers a smaller
marginal response.
An alternative explanation is dilution: if the treatment effect is zero or negative
for high-stratum tracts, including them in S2-pooled mechanically pulls the average
toward zero.
The thin Stratum~2+ control pool prevents us from distinguishing these channels.

Third, the \emph{credibility of causal inference} diminishes sharply once we
relax the first-fire restriction for high-stratum tracts.
Tracts with two or more prior fire years (Stratum~2+) have a WFP profile so extreme
that no adequate comparison group exists among never-treated tracts; the 49 controls
that pass the WFP floor are insufficient.
This is the structural limitation that motivates the baseline first-fire design:
it is not merely a sample-selection choice, but the boundary of the region where
credible propensity-score matching is feasible.

These results support treating the baseline finding as a lower bound on the effect
among all large-fire tracts, with the caveat that the first-fire design estimates
a specific estimand --- the ATT for tracts experiencing their first recorded large
wildfire --- that may not generalise to repeat-fire tracts.

\medskip
\noindent\textbf{Caution on Stratum~1.}
After the WFP floor and common support restriction, only 130 control tracts are
matched to 154 treated tracts in Stratum~1.
The pre-trend coefficient for Stratum~1 is $\hat{\beta}_{-2} = -0.56$ pp
(95\% CI: $[-1.32, +0.20]$), which is statistically indistinguishable from zero
but directionally negative.
Given the thin control pool, we interpret the Stratum~1 ATT ($+0.03$ pp) as imprecise
rather than as evidence of no effect, and we do not draw separate causal conclusions
from this stratum.

\medskip
\noindent\textbf{Caution on Stratum~2+.}
After applying the WFP floor ($\geq$40th percentile) and common support trimming,
only 49 control tracts remain in Stratum~2+.
This is insufficient for credible propensity-score inference.
We report event-study coefficients for Stratum~2+ in Figure~\ref{fig:het_strata}
for descriptive completeness, but we do not draw causal conclusions from that column.
Formal inference is restricted to the ATT$_1 -$ ATT$_0$ comparison.

\begin{table}[ht]
\centering
\caption{ATT by Prior Fire Stratum: In-migration Rate (pp)}
\label{tab:app_het_strata}
\small
\begin{tabular}{lcccccc}
\toprule
 & \multicolumn{2}{c}{Stratum 0} & \multicolumn{2}{c}{Stratum 1} & \multicolumn{2}{c}{Stratum 2+} \\
 & \multicolumn{2}{c}{(0 prior fire yrs)} & \multicolumn{2}{c}{(1 prior fire yr)} & \multicolumn{2}{c}{(2+ prior fire yrs)} \\
\cmidrule(lr){2-3}\cmidrule(lr){4-5}\cmidrule(lr){6-7}
 & ATT (pp) & 95\% CI & ATT (pp) & 95\% CI & ATT (pp) & 95\% CI \\
\midrule
In-migration rate    & $+0.65$ & $[-0.14,\; +1.44]$ & $+0.03$ & $[-1.40,\; +1.46]$ & $+0.19$ & $[-1.90,\; +2.28]$ \\
\midrule
$\beta_{-2}$ (pre-trend) & $-0.06$ & $[-0.60,\; +0.49]$ & $-0.56$ & $[-1.32,\; +0.20]$ & --- & --- \\
\midrule
$N$ treated & \multicolumn{2}{c}{$214$} & \multicolumn{2}{c}{$154$} & \multicolumn{2}{c}{$695$} \\
$N$ controls & \multicolumn{2}{c}{$8,\!139$} & \multicolumn{2}{c}{$130$} & \multicolumn{2}{c}{$49$$^\dagger$} \\
\bottomrule
\end{tabular}
\medskip

\noindent\footnotesize
\textit{Notes:} Stratum defined by count of calendar years with $\geq$1 large MTBS wildfire
($\geq$1,000 acres) in 1984--2014 overlapping the tract.
Control group: never-treated tracts in same stratum.
Stratum-specific PS-IPW; standard errors clustered by county.
Aggregate ATT $= (\hat{\beta}_0 + \hat{\beta}_{+1} + \hat{\beta}_{+2})/3$.
Heterogeneity test (bootstrap, $B = 200$, county block):
ATT$_1 -$ ATT$_0 =$ $-0.54$ pp, $p =$ $0.45$.\\
$^\dagger$ Stratum~2+ control pool is too thin for credible inference; reported for completeness only.
\end{table}

\begin{figure}[ht]
\centering
\includegraphics[width=\textwidth]{../results/fig_firecount_heterogeneity.png}
\caption{Event-Study DiD by Prior Fire Stratum: In-migration Rate}
\label{fig:het_strata}
\medskip

\noindent\footnotesize
\textit{Notes:} Three-panel event-study plot for Strata~0 (top), 1 (middle), and 2+ (bottom).
Each panel uses stratum-specific PS-IPW matching on WFP 2012 raster and 2014 ACS covariates.
Shaded bands: 95\% confidence intervals, cluster-SE by county.
Reference period $h = -1$ (ACS 2014) normalized to zero.
Triangles (Stratum~2+ panel) indicate estimates from a thin control pool ($N = 51$);
interpret with caution.
\end{figure}

\end{document}
